% Options for packages loaded elsewhere
\PassOptionsToPackage{unicode}{hyperref}
\PassOptionsToPackage{hyphens}{url}
\documentclass[
  spanish,
  11pt,
  letterpaper,
]{article}
\usepackage{xcolor}
\usepackage[margin=2.5cm]{geometry}
\usepackage{amsmath,amssymb}
\setcounter{secnumdepth}{-\maxdimen} % remove section numbering
\usepackage{iftex}
\ifPDFTeX
  \usepackage[T1]{fontenc}
  \usepackage[utf8]{inputenc}
  \usepackage{textcomp} % provide euro and other symbols
\else % if luatex or xetex
  \usepackage{unicode-math} % this also loads fontspec
  \defaultfontfeatures{Scale=MatchLowercase}
  \defaultfontfeatures[\rmfamily]{Ligatures=TeX,Scale=1}
\fi
\usepackage{lmodern}
\ifPDFTeX\else
  % xetex/luatex font selection
\fi
% Use upquote if available, for straight quotes in verbatim environments
\IfFileExists{upquote.sty}{\usepackage{upquote}}{}
\IfFileExists{microtype.sty}{% use microtype if available
  \usepackage[]{microtype}
  \UseMicrotypeSet[protrusion]{basicmath} % disable protrusion for tt fonts
}{}
\makeatletter
\@ifundefined{KOMAClassName}{% if non-KOMA class
  \IfFileExists{parskip.sty}{%
    \usepackage{parskip}
  }{% else
    \setlength{\parindent}{0pt}
    \setlength{\parskip}{6pt plus 2pt minus 1pt}}
}{% if KOMA class
  \KOMAoptions{parskip=half}}
\makeatother
\ifLuaTeX
\usepackage[bidi=basic,shorthands=off]{babel}
\else
\usepackage[bidi=default,shorthands=off]{babel}
\fi
\ifLuaTeX
  \usepackage{selnolig} % disable illegal ligatures
\fi
\setlength{\emergencystretch}{3em} % prevent overfull lines
\providecommand{\tightlist}{%
  \setlength{\itemsep}{0pt}\setlength{\parskip}{0pt}}
\usepackage{bookmark}
\IfFileExists{xurl.sty}{\usepackage{xurl}}{} % add URL line breaks if available
\urlstyle{same}
\hypersetup{
  pdflang={es},
  hidelinks,
  pdfcreator={LaTeX via pandoc}}

\author{}
\date{}

\begin{document}

Propuesta de Solución Tecnológica: Sistema de Órdenes de Servicio y
Citas

Proyecto: Automatización y Control del Ciclo de Servicio

Empresa Solicitante: Servicio Automotriz Hermanos Ruiz, S.A. de C.V.

Equipo Proveedor: Equipo B (5 integrantes)

\begin{enumerate}
\def\labelenumi{\arabic{enumi}.}
\item
  Resumen Ejecutivo y Enfoque
\item
\end{enumerate}

Nuestro objetivo es transformar la operación manual actual (cuadernos,
llamadas constantes y órdenes en papel) en una plataforma digital ágil,
transparente y segura. Trabajaremos bajo la metodología Scrum en sprints
de dos semanas, garantizando entregas incrementales y revisiones
continuas (\emph{Sprint Reviews}) para asegurar que el sistema se adapte
exactamente a las necesidades reales del taller.

\begin{enumerate}
\def\labelenumi{\arabic{enumi}.}
\setcounter{enumi}{2}
\item
  solución a los problemas actuales del taller
\end{enumerate}

La propuesta de software ataca directamente los 5 puntos críticos que
hoy le cuestan dinero y clientes a la empresa:

\begin{itemize}
\item
  \textbf{Fin a las disputas por costos}: Toda cotización y su
  respectiva autorización quedarán registradas digitalmente con fecha,
  hora, monto y medio de aprobación, eliminando el "dijo que sí, pero
  cobra otra cosa".
\item
  \textbf{Reducción de llamadas innecesarias:} Implementaremos un
  sistema de notificaciones y un enlace seguro web para que los clientes
  consulten el estado de su vehículo desde su celular sin saturar a la
  recepcionista.
\item
  \textbf{Control de vehículos y bahías}: Visualización en tiempo real
  de las 4 bahías de trabajo y el estatus de las refacciones para evitar
  unidades detenidas sin justificación.
\item
  \textbf{Historial vehicular unificado}: Cada vez que un cliente
  regrese, el sistema mostrará el historial completo de reparaciones
  anteriores y garantías vigentes.
\item
  \textbf{Evidencia contra reclamos por daños:} El módulo de recepción
  incluirá un \emph{checklist} digital y carga fotográfica obligatoria
  del estado inicial del vehículo.
\end{itemize}

\begin{enumerate}
\def\labelenumi{\arabic{enumi}.}
\setcounter{enumi}{3}
\item
  Propuesta de Arquitectura y Alcance Funcional
\end{enumerate}

Para cumplir con las restricciones del proyecto (uso exclusivo de
herramientas gratuitas, de código abierto y planes educativos),
proponemos una arquitectura moderna basada en microservicios y APIs:

\begin{itemize}
\item
  Backend / Lógica de Negocio: Desarrollado en Python con FastAPI,
  garantizando alta velocidad de respuesta, seguridad y la generación
  automática de documentación técnica.
\item
  Interfaz de Usuario (Frontend):

  \begin{itemize}
  \item
    \emph{Para el personal (Gerente, Recepción, Mecánicos, Almacén):}
    Una interfaz web optimizada (con botones grandes y accesibles para
    tabletas en el área de trabajo con ruido y manos sucias).
  \item
    \emph{Para el Cliente:} Una vista web responsiva accesible mediante
    un enlace seguro enviado por correo o mensaje, sin necesidad de
    instalar aplicaciones móviles.
  \end{itemize}
\item
  Control de Versiones y Despliegue: Repositorio centralizado en GitHub
  de la organización del curso, asegurando transparencia y trazabilidad
  mediante \emph{Pull Requests} y ramas de desarrollo.
\end{itemize}

\begin{enumerate}
\def\labelenumi{\arabic{enumi}.}
\setcounter{enumi}{4}
\item
  Enfoque del Primer Producto Mínimo Funcional (Release v0.1.0)
\end{enumerate}

Para la Fase 1, nos enfocaremos en construir y desplegar el flujo núcleo
operativo de extremo a extremo:

\begin{enumerate}
\def\labelenumi{\arabic{enumi}.}
\item
  Recepción: La recepcionista registra al cliente, vehículo,
  kilometraje, nivel de combustible y fotos de ingreso.
\item
  Diagnóstico y Asignación: La orden se asigna a un mecánico y a una de
  las 4 bahías disponibles, registrando el diagnóstico técnico.
\item
  Cotización y Autorización: Se genera el presupuesto de mano de obra y
  refacciones, enviándolo al cliente para su validación digital
  registrada.
\item
  Monitoreo Gerencial: El gerente cuenta con un panel para visualizar
  las órdenes abiertas, sus estados actuales y la carga de trabajo del
  taller.
\end{enumerate}

\textbf{5. Consideraciones Legales y de Seguridad}

Debido al manejo de información de clientes, el sistema cumplirá
estrictamente con las directrices de la Ley Federal de Protección de
Datos Personales en Posesión de los Particulares (DOF), incorporando
avisos de privacidad y uso exclusivo de datos ficticios durante las
fases de prueba y desarrollo.

Agregar historia, requerimientos y la entrevista al cliente.

\textbf{6. Requerimientos Extraídos del Anteproyecto}

A partir del documento del anteproyecto, podemos clasificar los
requerimientos iniciales de la siguiente manera:

\begin{itemize}
\item
  Requerimientos Funcionales (Lo que el sistema debe hacer):

  \begin{enumerate}
  \def\labelenumi{\arabic{enumi}.}
  \item
    El sistema debe permitir el registro y búsqueda de clientes y sus
    vehículos (placas, marca, modelo, año).
  \item
    El sistema debe permitir abrir una orden de servicio registrando
    kilometraje, nivel de combustible, checklist y fotos de daños
    previos.
  \item
    El sistema debe permitir la asignación de órdenes de trabajo a
    mecánicos y a una de las 4 bahías del taller.
  \item
    El sistema debe generar cotizaciones de mano de obra y refacciones,
    enviándolas al cliente para su autorización digital.
  \item
    El sistema debe permitir al cliente consultar el estado de su
    vehículo mediante un enlace seguro desde su celular sin instalar
    aplicaciones.
  \item
    El gerente debe contar con un panel de control para ver las órdenes
    abiertas, tiempos, ingresos y carga de trabajo de las bahías.
  \end{enumerate}
\item
  Requerimientos No Funcionales / Restricciones (Cómo debe comportarse):

  \begin{enumerate}
  \def\labelenumi{\arabic{enumi}.}
  \item
    Restricción de Presupuesto: No se deben pagar licencias en esta
    fase; se deben usar herramientas gratuitas, de código abierto o
    planes educativos.
  \item
    Usabilidad: Los mecánicos necesitan una interfaz con botones
    grandes, pocos pasos y optimizada para uso en tabletas o celulares
    (con manos sucias y ruido en el área de trabajo).
  \item
    Seguridad y Trazabilidad: La autorización del cliente debe registrar
    fecha, hora, medio y monto de forma inalterable.
  \item
    Cumplimiento Legal: El sistema debe cumplir con la Ley Federal de
    Protección de Datos Personales (aviso de privacidad y medidas de
    seguridad), utilizando datos ficticios para pruebas.
  \end{enumerate}
\end{itemize}

7. Preguntas Estratégicas para el Negocio

Para afinar el alcance y entender la viabilidad técnica y operativa,
estas son las preguntas clave que le pueden hacer al cliente (Gerencia):

\begin{itemize}
\item
  Sobre dispositivos y recursos económicos:

  \begin{itemize}
  \item
    \emph{"Dado que los mecánicos trabajarán en el área de servicio y
    requerimos una interfaz con tabletas o celulares, ¿el taller ya
    cuenta con dispositivos móviles (tabletas) para el personal o el
    equipo proveedor debe presupuestar/considerar el uso de los
    teléfonos personales de los mecánicos (BYOD)?"}
  \item
    \emph{"Para el envío de notificaciones y cotizaciones a los clientes
    (vía correo electrónico o posibles futuros canales), ¿cuentan con un
    servidor de correo corporativo o prefieren que utilicemos servicios
    gratuitos en la nube (como SMTP de pruebas o APIs básicas)?"}
  \end{itemize}
\item
  Sobre las reglas del negocio y el flujo operativo:

  \begin{itemize}
  \item
    "Mencionan que una cotización tiene vigencia de 5 días naturales y
    que si un trabajo adicional incrementa el costo en más del 10\%
    requiere reautorización. ¿El sistema debe bloquear automáticamente
    la orden hasta que llegue esa nueva autorización o solo debe alertar
    al asesor?"
  \item
    "Para el caso de las 3 pequeñas flotillas empresariales que
    atienden, ¿el crédito autorizado lo maneja directamente el gerente
    con aprobación manual o el sistema debe registrar un límite de
    crédito por empresa?"
  \end{itemize}
\item
  Sobre la recepción y evidencias:

  \begin{itemize}
  \item
    "Para las fotografías de recepción de los vehículos, ¿existe alguna
    restricción de almacenamiento o tamaño máximo de las imágenes para
    evitar saturar el servidor gratuito que utilizaremos en el
    proyecto?"
  \end{itemize}
\end{itemize}

8. Ejemplo de Historia de Usuario basada en el Anteproyecto

Tomando un requisito crítico del documento (la autorización de la
cotización), así se estructura una historia de usuario lista para el
Product Backlog:

\begin{itemize}
\item
  Título: Autorización digital de cotización por parte del cliente.
\item
  Historia de Usuario:
\end{itemize}

Como cliente del taller, quiero recibir un enlace seguro en mi celular
para revisar el detalle de la cotización (mano de obra y refacciones) y
poder aceptarla o rechazarla con un clic, para autorizar los trabajos de
mi vehículo de forma rápida y sin necesidad de llamadas telefónicas o
visitas presenciales.

\begin{itemize}
\item
  Criterios de Aceptación:

  \begin{enumerate}
  \def\labelenumi{\arabic{enumi}.}
  \item
    El sistema genera un enlace web único y seguro al crear la
    cotización.
  \item
    Al abrir el enlace en cualquier navegador móvil, el cliente
    visualiza el desglose de la cotización (vigente por 5 días).
  \item
    La interfaz cuenta con dos botones claros: \emph{"Autorizar"} y
    \emph{"Rechazar"}.
  \item
    Al hacer clic en autorizar, el sistema registra de forma inalterable
    la fecha, hora, monto y dirección IP, cambiando el estatus de la
    orden de servicio automáticamente.
  \end{enumerate}
\end{itemize}

\end{document}
